\documentclass{article}\usepackage[]{graphicx}\usepackage[]{xcolor}
% maxwidth is the original width if it is less than linewidth
% otherwise use linewidth (to make sure the graphics do not exceed the margin)
\makeatletter
\def\maxwidth{ %
  \ifdim\Gin@nat@width>\linewidth
    \linewidth
  \else
    \Gin@nat@width
  \fi
}
\makeatother

\definecolor{fgcolor}{rgb}{0.345, 0.345, 0.345}
\newcommand{\hlnum}[1]{\textcolor[rgb]{0.686,0.059,0.569}{#1}}%
\newcommand{\hlsng}[1]{\textcolor[rgb]{0.192,0.494,0.8}{#1}}%
\newcommand{\hlcom}[1]{\textcolor[rgb]{0.678,0.584,0.686}{\textit{#1}}}%
\newcommand{\hlopt}[1]{\textcolor[rgb]{0,0,0}{#1}}%
\newcommand{\hldef}[1]{\textcolor[rgb]{0.345,0.345,0.345}{#1}}%
\newcommand{\hlkwa}[1]{\textcolor[rgb]{0.161,0.373,0.58}{\textbf{#1}}}%
\newcommand{\hlkwb}[1]{\textcolor[rgb]{0.69,0.353,0.396}{#1}}%
\newcommand{\hlkwc}[1]{\textcolor[rgb]{0.333,0.667,0.333}{#1}}%
\newcommand{\hlkwd}[1]{\textcolor[rgb]{0.737,0.353,0.396}{\textbf{#1}}}%
\let\hlipl\hlkwb

\usepackage{framed}
\makeatletter
\newenvironment{kframe}{%
 \def\at@end@of@kframe{}%
 \ifinner\ifhmode%
  \def\at@end@of@kframe{\end{minipage}}%
  \begin{minipage}{\columnwidth}%
 \fi\fi%
 \def\FrameCommand##1{\hskip\@totalleftmargin \hskip-\fboxsep
 \colorbox{shadecolor}{##1}\hskip-\fboxsep
     % There is no \\@totalrightmargin, so:
     \hskip-\linewidth \hskip-\@totalleftmargin \hskip\columnwidth}%
 \MakeFramed {\advance\hsize-\width
   \@totalleftmargin\z@ \linewidth\hsize
   \@setminipage}}%
 {\par\unskip\endMakeFramed%
 \at@end@of@kframe}
\makeatother

\definecolor{shadecolor}{rgb}{.97, .97, .97}
\definecolor{messagecolor}{rgb}{0, 0, 0}
\definecolor{warningcolor}{rgb}{1, 0, 1}
\definecolor{errorcolor}{rgb}{1, 0, 0}
\newenvironment{knitrout}{}{} % an empty environment to be redefined in TeX

\usepackage{alltt}
\usepackage[margin=1.0in]{geometry} % To set margins
\usepackage{amsmath}  % This allows me to use the align functionality.
                      % If you find yourself trying to replicate
                      % something you found online, ensure you're
                      % loading the necessary packages!
\usepackage{amsfonts} % Math font
\usepackage{fancyvrb}
\usepackage{hyperref} % For including hyperlinks
\usepackage[shortlabels]{enumitem}% For enumerated lists with labels specified
                    % We had to run tlmgr_install("enumitem") in R
\usepackage{float}    % For telling R where to put a table/figure
\usepackage{natbib}        %For the bibliography
\usepackage[skins,breakable]{tcolorbox}        %For the bibliography
\bibliographystyle{apalike}%For the bibliography
\IfFileExists{upquote.sty}{\usepackage{upquote}}{}
\begin{document}
\noindent\textbf{Homework 6: Regression}\\
\noindent\makebox[\linewidth]{\rule{\textwidth}{0.4pt}}
\noindent\underline{Submission Notes:} Please write your answers to each 
question under that question (not altogether at the end). It will be easier for 
me to give full credit if your answer is commented well, including comments 
about where you are completing each sub-part of the question. For an example,
see Homework 1.\\
\noindent\makebox[\linewidth]{\rule{\textwidth}{0.4pt}}


\begin{enumerate}
  %%%%%%%%%%%%%%%%%%%%%%%%%%%%%%%%%%%%%%%%%%%%%%%%%%%%%%%%%%%%%%%%%%%%%%%%%%%%%%
  %%%%%%%%%%%%%%%%%%%%%%%%%%%%%%%%%%%%%%%%%%%%%%%%%%%%%%%%%%%%%%%%%%%%%%%%%%%%%%
  % Question 1
  %%%%%%%%%%%%%%%%%%%%%%%%%%%%%%%%%%%%%%%%%%%%%%%%%%%%%%%%%%%%%%%%%%%%%%%%%%%%%%
  %%%%%%%%%%%%%%%%%%%%%%%%%%%%%%%%%%%%%%%%%%%%%%%%%%%%%%%%%%%%%%%%%%%%%%%%%%%%%%
  \item In Lecture 13, we discussed the linear algebra formulation of linear regression.
  For example, we can set up perfectly simulated regression data as follows:
\begin{knitrout}\scriptsize
\definecolor{shadecolor}{rgb}{0.969, 0.969, 0.969}\color{fgcolor}\begin{kframe}
\begin{alltt}
\hldef{simulate.reg} \hlkwb{<-} \hlkwa{function}\hldef{(}\hlkwc{n}\hldef{=}\hlnum{100}\hldef{,}
                         \hlkwc{betas}\hldef{=}\hlkwd{c}\hldef{(}\hlnum{4.35}\hldef{,} \hlnum{1.73}\hldef{,} \hlnum{0.02}\hldef{,} \hlnum{2.78}\hldef{),}
                         \hlkwc{sigma}\hldef{=}\hlnum{4.16}\hldef{,}
                         \hlkwc{set.seed}\hldef{=F)\{}
  \hlcom{# Specify Coefficients}
  \hldef{beta} \hlkwb{<-} \hlkwd{matrix}\hldef{(}\hlkwc{data} \hldef{= betas,} \hlkwc{nrow}\hldef{=}\hlnum{4}\hldef{)}
  \hlcom{# Specify Error Distribution}
  \hldef{mu.epsilon} \hlkwb{<-} \hlnum{0}
  \hldef{sigma.epsilon}\hlkwb{<-}\hlkwd{sqrt}\hldef{(sigma)}
  \hlcom{# Simulate Regression Data}
  \hlkwa{if}\hldef{(set.seed)\{}\hlkwd{set.seed}\hldef{(}\hlnum{7272}\hldef{)\}}
  \hldef{x1} \hlkwb{<-} \hlkwd{runif}\hldef{(}\hlkwc{n}\hldef{=n,} \hlkwc{min}\hldef{=}\hlopt{-}\hlnum{5}\hldef{,} \hlkwc{max}\hldef{=}\hlnum{5}\hldef{)}
  \hldef{x2} \hlkwb{<-} \hlkwd{runif}\hldef{(}\hlkwc{n}\hldef{=n,} \hlkwc{min}\hldef{=}\hlopt{-}\hlnum{5}\hldef{,} \hlkwc{max}\hldef{=}\hlnum{5}\hldef{)}
  \hldef{x} \hlkwb{<-} \hlkwd{cbind}\hldef{(}\hlnum{1}\hldef{, x1, x2, x1}\hlopt{*}\hldef{x2)}
  \hldef{epsilon} \hlkwb{<-} \hlkwd{rnorm}\hldef{(}\hlkwc{n}\hldef{=n,} \hlkwc{mean}\hldef{=}\hlnum{0}\hldef{,} \hlkwc{sd}\hldef{=sigma.epsilon)}
  \hldef{y} \hlkwb{<-} \hldef{x} \hlopt{%*%} \hldef{beta} \hlopt{+} \hldef{epsilon}
  \hlcom{# Return simulated data}
  \hldef{dat} \hlkwb{<-} \hlkwd{tibble}\hldef{(}\hlkwc{y}\hldef{=y,} \hlkwc{x1}\hldef{=x1,} \hlkwc{x2}\hldef{=x2)}
  \hlkwd{return}\hldef{(dat)}
\hldef{\}}
\end{alltt}
\end{kframe}
\end{knitrout}
  In Lecture 15, we discussed regression-based inference. During office hours, 
  I had some discussions about the mathematical statistics behind these results. 
  Although we don't do the real-analysis style derivations in this course that 
  we might in Mathematical Statistics, we have plenty of tools to investigate 
  these ideas numerically. Complete the parts below to conduct this exploration.
  \begin{enumerate}
    \item Use the \texttt{simulate.reg()} function to simulate regression data 
    using the default set of coefficients over a grid of $\sigma_\epsilon$ 
    values from $\sigma_\epsilon=0$ to $\sigma_\epsilon=5$. Compute the
    $SS_{res}$ and $SS_{reg}$ from the resulting regression model, and plot
    the points using color to denote the $\sigma_\epsilon$ value. Make sure you
    set \texttt{set.seed} to \texttt{T} so you can compare the output at
    fixed $\mathbf{x}$ and $\mathbf{Y}$. \textbf{Why do you see the pattern
    that emerges?}\\
    \textbf{Hint:} I found it was easier to manually compute $SS_{res}$ and
    $SS_{reg}$ than it was to extract them from the result of an \texttt{aov}
    call.
    
    \begin{figure}[H]
    \centering
    
\begin{knitrout}\scriptsize
\definecolor{shadecolor}{rgb}{0.969, 0.969, 0.969}\color{fgcolor}\begin{kframe}
\begin{alltt}
\hldef{asim} \hlkwb{<-} \hlkwd{tibble}\hldef{(}\hlkwc{sigma} \hldef{=} \hlkwd{rep}\hldef{(}\hlnum{NA}\hldef{,} \hlnum{1000}\hldef{),}
           \hlkwc{SSres} \hldef{=} \hlkwd{rep}\hldef{(}\hlnum{NA}\hldef{,} \hlnum{1000}\hldef{),}
           \hlkwc{SSreg} \hldef{=} \hlkwd{rep}\hldef{(}\hlnum{NA}\hldef{,} \hlnum{1000}\hldef{))}

\hldef{sigmas} \hlkwb{<-} \hlkwd{seq}\hldef{(}\hlkwc{from} \hldef{=} \hlnum{0}\hldef{,} \hlkwc{to} \hldef{=} \hlnum{5}\hldef{,} \hlkwc{length.out} \hldef{=} \hlnum{1000}\hldef{)}

\hlkwa{for}\hldef{(i} \hlkwa{in} \hlnum{1}\hlopt{:}\hlnum{1000}\hldef{)\{}

  \hldef{dat} \hlkwb{<-} \hlkwd{simulate.reg}\hldef{(}\hlkwc{sigma} \hldef{= sigmas[i],} \hlkwc{set.seed}\hldef{=T)}

  \hldef{asim}\hlopt{$}\hldef{sigma[i]} \hlkwb{<-} \hldef{sigmas[i]}

  \hldef{anova} \hlkwb{<-} \hlkwd{aov}\hldef{(y} \hlopt{~} \hldef{x1} \hlopt{+} \hldef{x2} \hlopt{+} \hldef{x1}\hlopt{*}\hldef{x2,} \hlkwc{data} \hldef{= dat)}

  \hldef{asim}\hlopt{$}\hldef{SSres[i]} \hlkwb{<-} \hlkwd{sum}\hldef{((dat}\hlopt{$}\hldef{y} \hlopt{-} \hldef{anova}\hlopt{$}\hldef{residuals)}\hlopt{^}\hlnum{2}\hldef{)}

  \hldef{asim}\hlopt{$}\hldef{SSreg[i]} \hlkwb{<-} \hlkwd{sum}\hldef{((anova}\hlopt{$}\hldef{residuals} \hlopt{-} \hlkwd{mean}\hldef{(dat}\hlopt{$}\hldef{y))}\hlopt{^}\hlnum{2}\hldef{)}

\hldef{\}}

\hlkwd{ggplot}\hldef{(}\hlkwc{data} \hldef{= asim,} \hlkwd{aes}\hldef{(}\hlkwc{x} \hldef{= SSres,} \hlkwc{y} \hldef{= SSreg,} \hlkwc{color} \hldef{= sigma))} \hlopt{+}
  \hlkwd{geom_point}\hldef{()} \hlopt{+}
  \hlkwd{xlab}\hldef{(}\hlkwd{bquote}\hldef{(SS[res]))} \hlopt{+}
  \hlkwd{ylab}\hldef{(}\hlkwd{bquote}\hldef{(SS[reg]))} \hlopt{+}
  \hlkwd{labs}\hldef{(}\hlkwc{color} \hldef{=} \hlkwd{bquote}\hldef{(sigma[epsilon]))} \hlopt{+}
  \hlkwd{theme_bw}\hldef{()}
\end{alltt}
\end{kframe}
\includegraphics[width=\maxwidth]{figure/unnamed-chunk-3-1} 
\end{knitrout}
    
    
    \label{sigma.fig}
    \caption{$\sigma$ across levels of $SS_{res}$ and $SS_{reg}$ (fixed $\mathbf{x}$ and $\mathbf{Y}$)}
    \end{figure}
    
    \textcolor{red}{We see that higher values of $SS_{res}$ and $SS_{reg}$ are associated with higher values of $\sigma_{\epsilon}$. This makes sense as $\sigma_{\epsilon}$ affects the variation around how the data fits the model, and higher variation in $\sigma_{\epsilon}$ will result in higher levels of variation in the regression residuals and against $\bar{Y}$.}
    
    \item Use the \texttt{simulate.reg()} function to simulate regression data 
    using the default set of coefficients. Create a plot that shows the 
    distribution of the $\beta$s. Make sure you set \texttt{set.seed} to 
    \texttt{F} so you can compare the output for differing samples of
    $\mathbf{x}$ and $\mathbf{Y}$. \textbf{What do you notice about their 
    distributions?}
    
    \begin{figure}[H]
    \centering
    
\begin{knitrout}\scriptsize
\definecolor{shadecolor}{rgb}{0.969, 0.969, 0.969}\color{fgcolor}\begin{kframe}
\begin{alltt}
\hldef{bsim} \hlkwb{<-} \hlkwd{tibble}\hldef{(}\hlkwc{Beta0} \hldef{=} \hlkwd{rep}\hldef{(}\hlnum{NA}\hldef{,} \hlnum{1000}\hldef{),}
           \hlkwc{Beta1} \hldef{=} \hlkwd{rep}\hldef{(}\hlnum{NA}\hldef{,} \hlnum{1000}\hldef{),}
           \hlkwc{Beta2} \hldef{=} \hlkwd{rep}\hldef{(}\hlnum{NA}\hldef{,} \hlnum{1000}\hldef{),}
           \hlkwc{Beta3} \hldef{=} \hlkwd{rep}\hldef{(}\hlnum{NA}\hldef{,} \hlnum{1000}\hldef{))}

\hlkwa{for} \hldef{(i} \hlkwa{in} \hlnum{1}\hlopt{:}\hlnum{1000}\hldef{)\{}

  \hldef{dat} \hlkwb{<-} \hlkwd{simulate.reg}\hldef{(}\hlkwc{set.seed}\hldef{=F)}

  \hldef{mod} \hlkwb{<-} \hlkwd{data.frame}\hldef{(}\hlkwd{summary}\hldef{(}\hlkwd{lm}\hldef{(y} \hlopt{~} \hldef{x1} \hlopt{+} \hldef{x2} \hlopt{+} \hldef{x1}\hlopt{*}\hldef{x2,} \hlkwc{data} \hldef{= dat))}\hlopt{$}\hldef{coefficients)}

  \hldef{bsim}\hlopt{$}\hldef{Beta0[i]} \hlkwb{<-} \hldef{mod}\hlopt{$}\hldef{Estimate[}\hlnum{1}\hldef{]}
  \hldef{bsim}\hlopt{$}\hldef{Beta1[i]} \hlkwb{<-} \hldef{mod}\hlopt{$}\hldef{Estimate[}\hlnum{2}\hldef{]}
  \hldef{bsim}\hlopt{$}\hldef{Beta2[i]} \hlkwb{<-} \hldef{mod}\hlopt{$}\hldef{Estimate[}\hlnum{3}\hldef{]}
  \hldef{bsim}\hlopt{$}\hldef{Beta3[i]} \hlkwb{<-} \hldef{mod}\hlopt{$}\hldef{Estimate[}\hlnum{4}\hldef{]}


\hldef{\}}

\hldef{bsim} \hlkwb{<-} \hldef{bsim |>}
  \hlkwd{pivot_longer}\hldef{(}\hlkwc{cols} \hldef{=} \hlkwd{colnames}\hldef{(bsim),}
           \hlkwc{names_to} \hldef{=} \hlsng{"Beta"}\hldef{,}
           \hlkwc{values_to} \hldef{=} \hlsng{"Value"}\hldef{)}


\hlkwd{library}\hldef{(ggridges)}
\hlkwd{ggplot}\hldef{(}\hlkwc{data} \hldef{= bsim,} \hlkwd{aes}\hldef{(}\hlkwc{x} \hldef{= Value,} \hlkwc{y} \hldef{=} \hlkwd{fct_rev}\hldef{(}\hlkwd{as.factor}\hldef{(Beta))))} \hlopt{+}
  \hlkwd{geom_density_ridges}\hldef{(}\hlkwc{alpha} \hldef{=} \hlnum{0.5}\hldef{)} \hlopt{+}
  \hlkwd{theme_bw}\hldef{()} \hlopt{+}
  \hlkwd{xlab}\hldef{(}\hlsng{"Estimate"}\hldef{)} \hlopt{+}
  \hlkwd{geom_hline}\hldef{(}\hlkwc{yintercept} \hldef{=} \hlnum{0}\hldef{)} \hlopt{+}
  \hlkwd{ylab}\hldef{(}\hlsng{''}\hldef{)} \hlopt{+}
  \hlkwd{scale_y_discrete}\hldef{(}
\hlkwc{labels} \hldef{=} \hlkwd{c}\hldef{(}
  \hlsng{"Beta0"} \hldef{=} \hlkwd{bquote}\hldef{(beta[}\hlnum{0}\hldef{]),}
  \hlsng{"Beta1"} \hldef{=} \hlkwd{bquote}\hldef{(beta[}\hlnum{1}\hldef{]),}
  \hlsng{"Beta2"} \hldef{=} \hlkwd{bquote}\hldef{(beta[}\hlnum{2}\hldef{]),}
  \hlsng{"Beta3"} \hldef{=} \hlkwd{bquote}\hldef{(beta[}\hlnum{3}\hldef{])}
\hldef{)}
  \hldef{)}
\end{alltt}
\end{kframe}
\includegraphics[width=\maxwidth]{figure/unnamed-chunk-4-1} 
\end{knitrout}
    
    \label{beta.dist}
    \caption{Distribution of the $\beta$ coefficients}
    
    \end{figure}
    
    \textcolor{red}{The spread of their distributions is very different. $\beta_{0}$ has the widest spread, followed by non-interaction terms $\beta_{1}$ and $\beta_{2}$. Then, the coefficient with the least spread is clearly the interaction term $\beta_{3}$}
    
    \item Use the \texttt{simulate.reg()} function to simulate regression data 
    using the default set of coefficients. Create a plot that shows the 
    distribution of the $\hat{y}$s at $x_1=3.6$ and $x_2=-1.9$. Make sure you 
    set \texttt{set.seed} to \texttt{F} so you can compare the output for 
    differing samples of $\mathbf{x}$ and $\mathbf{Y}$. \textbf{What do you 
    notice about the distribution?}
    
    \begin{figure}[H]
    \centering
    
\begin{knitrout}\scriptsize
\definecolor{shadecolor}{rgb}{0.969, 0.969, 0.969}\color{fgcolor}\begin{kframe}
\begin{alltt}
\hldef{csim} \hlkwb{<-} \hlkwd{tibble}\hldef{(}\hlkwc{yhat} \hldef{=} \hlkwd{rep}\hldef{(}\hlnum{NA}\hldef{,} \hlnum{1000}\hldef{))}

\hlkwa{for} \hldef{(i} \hlkwa{in} \hlnum{1}\hlopt{:}\hlnum{1000}\hldef{)\{}


  \hldef{dat} \hlkwb{<-} \hlkwd{simulate.reg}\hldef{(}\hlkwc{set.seed}\hldef{=F)}

  \hldef{mod} \hlkwb{<-} \hlkwd{data.frame}\hldef{(}\hlkwd{summary}\hldef{(}\hlkwd{lm}\hldef{(y} \hlopt{~} \hldef{x1} \hlopt{+} \hldef{x2} \hlopt{+} \hldef{x1}\hlopt{*}\hldef{x2,} \hlkwc{data} \hldef{= dat))}\hlopt{$}\hldef{coefficients)}

  \hldef{betas} \hlkwb{<-} \hldef{mod}\hlopt{$}\hldef{Estimate}

  \hldef{csim}\hlopt{$}\hldef{yhat[i]} \hlkwb{<-} \hldef{betas[}\hlnum{1}\hldef{]} \hlopt{+} \hlnum{3.6}\hlopt{*}\hldef{betas[}\hlnum{2}\hldef{]} \hlopt{-} \hlnum{1.9}\hlopt{*}\hldef{betas[}\hlnum{3}\hldef{]} \hlopt{-} \hlnum{3.6}\hlopt{*}\hlnum{1.9}\hlopt{*}\hldef{betas[}\hlnum{4}\hldef{]}

\hldef{\}}

\hldef{hist} \hlkwb{<-} \hlkwd{ggplot}\hldef{(}\hlkwc{data} \hldef{= csim,} \hlkwd{aes}\hldef{(}\hlkwc{x} \hldef{= yhat,} \hlkwc{y} \hldef{=} \hlkwd{after_stat}\hldef{(density)))} \hlopt{+}
  \hlkwd{geom_histogram}\hldef{(}\hlkwc{bins} \hldef{=} \hlnum{15}\hldef{,} \hlkwc{color} \hldef{=} \hlsng{"black"}\hldef{,} \hlkwc{fill} \hldef{=} \hlsng{"lightgrey"}\hldef{)} \hlopt{+}
  \hlkwd{theme_bw}\hldef{()} \hlopt{+}
  \hlkwd{geom_hline}\hldef{(}\hlkwc{yintercept} \hldef{=} \hlnum{0}\hldef{)} \hlopt{+}
  \hlkwd{xlab}\hldef{(}\hlkwd{bquote}\hldef{(}\hlkwd{hat}\hldef{(y)))} \hlopt{+}
  \hlkwd{scale_x_continuous}\hldef{(}\hlkwc{breaks} \hldef{=} \hlkwd{c}\hldef{(}\hlopt{-}\hlnum{10}\hldef{,} \hlopt{-}\hlnum{9}\hldef{,} \hlkwd{round}\hldef{(}\hlkwd{median}\hldef{(csim}\hlopt{$}\hldef{yhat),}\hlnum{2}\hldef{),} \hlopt{-}\hlnum{8}\hldef{,} \hlopt{-}\hlnum{7}\hldef{),} \hlkwc{limits} \hldef{=} \hlkwd{c}\hldef{(}\hlopt{-}\hlnum{10}\hldef{,}\hlopt{-}\hlnum{7}\hldef{))}

\hldef{boxplot} \hlkwb{<-} \hlkwd{ggplot}\hldef{(}\hlkwc{data} \hldef{= csim,} \hlkwd{aes}\hldef{(}\hlkwc{x} \hldef{=} \hlsng{""}\hldef{,} \hlkwc{y} \hldef{= yhat))} \hlopt{+}
  \hlkwd{geom_boxplot}\hldef{()} \hlopt{+}
  \hlkwd{coord_flip}\hldef{()} \hlopt{+}
  \hlkwd{xlab}\hldef{(}\hlsng{''}\hldef{)} \hlopt{+}
  \hlkwd{ylab}\hldef{(}\hlkwd{bquote}\hldef{(}\hlkwd{hat}\hldef{(y)))} \hlopt{+}
  \hlkwd{theme_bw}\hldef{()} \hlopt{+}
  \hlkwd{scale_y_continuous}\hldef{(}\hlkwc{breaks} \hldef{=} \hlkwd{c}\hldef{(}\hlopt{-}\hlnum{10}\hldef{,} \hlopt{-}\hlnum{9}\hldef{,} \hlkwd{round}\hldef{(}\hlkwd{median}\hldef{(csim}\hlopt{$}\hldef{yhat),} \hlnum{2}\hldef{),} \hlopt{-}\hlnum{8}\hldef{,} \hlopt{-}\hlnum{7}\hldef{),} \hlkwc{limits} \hldef{=} \hlkwd{c}\hldef{(}\hlopt{-}\hlnum{10}\hldef{,}\hlopt{-}\hlnum{7}\hldef{))}

\hlkwd{library}\hldef{(patchwork)}
\hldef{hist} \hlopt{+} \hldef{boxplot} \hlopt{+}
  \hlkwd{plot_layout}\hldef{(}\hlkwc{nrow} \hldef{=} \hlnum{2}\hldef{,} \hlkwc{heights} \hldef{=} \hlkwd{c}\hldef{(}\hlnum{3}\hldef{,} \hlnum{1}\hldef{))}
\end{alltt}
\end{kframe}
\includegraphics[width=\maxwidth]{figure/unnamed-chunk-5-1} 
\end{knitrout}
    
    \label{predict.dist}
    \caption{Distribution of $\hat{y}$'s at $x_{1} = 3.6$ and $x_{2} = -1.9$}
    
    \end{figure}
    
    \textcolor{red}{The distribution of the $\hat{y}$ is Gaussian and centered at -8.5. The IQR is 0.53 [-8.74, -8.21] which is similar to the initial parameter $\sigma$.}
    
    \item Describe why using $t$-based inference makes sense when conducting 
    hypothesis tests or computing confidence intervals for a coefficient $\beta_j$, 
    and the conditional expectation or prediction of $Y$?
  \end{enumerate}
  %%%%%%%%%%%%%%%%%%%%%%%%%%%%%%%%%%%%%%%%%%%%%%%%%%%%%%%%%%%%%%%%%%%%%%%%%%%%%%
  %%%%%%%%%%%%%%%%%%%%%%%%%%%%%%%%%%%%%%%%%%%%%%%%%%%%%%%%%%%%%%%%%%%%%%%%%%%%%%
  % Question 2
  %%%%%%%%%%%%%%%%%%%%%%%%%%%%%%%%%%%%%%%%%%%%%%%%%%%%%%%%%%%%%%%%%%%%%%%%%%%%%%
  %%%%%%%%%%%%%%%%%%%%%%%%%%%%%%%%%%%%%%%%%%%%%%%%%%%%%%%%%%%%%%%%%%%%%%%%%%%%%%
  \item \cite{Su25} investigate how well-being of preschool teachers influences 
  the effort they put into managing their emotions at work (emotional labor).
  They found that how committed a teacher is to their career (career commitment)
  plays a key role in linking well-being to emotional labor, and this connection 
  gets stronger when they have more social support from colleagues or the community.\\

  As part of this work, they aimed to replicate findings that the level of social support moderates the impact of 
  preschool teachers' career commitment on emotional labor. Below, we will assess this piece to the puzzle.\\

  Fit a regression model predicting \texttt{Emotional Labor} using 
  \texttt{career commitment}, \texttt{social support}, and their interaction.
  Also include Gender, Age, and teaching experience.
  
  \begin{enumerate}
    \item Perform a residual analysis. Do you have any concerns about the model
    fit?
    \item Perform an analysis of leverage/influence. Do you have any concerns 
    about the model fit?
    \item Provide the results of the overall $F$ test. What does this tell you
    about the model you fit?
    \item Based on the model fit, does the amount of social support impact the 
    effect of career commitment on emotional labor?
  \end{enumerate}
  %%%%%%%%%%%%%%%%%%%%%%%%%%%%%%%%%%%%%%%%%%%%%%%%%%%%%%%%%%%%%%%%%%%%%%%%%%%%%%
  %%%%%%%%%%%%%%%%%%%%%%%%%%%%%%%%%%%%%%%%%%%%%%%%%%%%%%%%%%%%%%%%%%%%%%%%%%%%%%
  % Question 3
  %%%%%%%%%%%%%%%%%%%%%%%%%%%%%%%%%%%%%%%%%%%%%%%%%%%%%%%%%%%%%%%%%%%%%%%%%%%%%%
  %%%%%%%%%%%%%%%%%%%%%%%%%%%%%%%%%%%%%%%%%%%%%%%%%%%%%%%%%%%%%%%%%%%%%%%%%%%%%%
  \item Continue with the regression from Question 2.
  \begin{enumerate}
    \item Use \texttt{emmeans()} from the \texttt{emmeans} package
    for \texttt{R} \citep{emmeans} to compute the estimated marginal means of
    \texttt{Emotional Labor} for low, average and high \texttt{career commitment}
    and \texttt{social support}. Report the results in a table and a figure.
    \item Use \texttt{emtrends()} from the \texttt{emmeans} package
    for \texttt{R} \citep{emmeans} to compute the slope for \texttt{career commitment}
    at low, average, and high \texttt{social support}. Report the results in a 
    table and use \texttt{ggeffect()} from the \texttt{ggeffects} package for
    \texttt{R} \citep{ggeffects} to create a plot. 
    \item Use \texttt{pairs()} to compare the slopes for \texttt{career commitment}
    at low, average, and high \texttt{social support}. What does this tell us?
    \item Use \texttt{johnson\_neyman()} from the \texttt{interactions} package
    for \texttt{R} \citep{interactions} to evaluate which values of 
    \text{social suppport} yield slopes for  \texttt{career commitment} that
    are statistically discernible from zero. What does this tell us?
  \end{enumerate}
\end{enumerate}
\bibliography{bib}
\end{document}
